\documentclass{amsart}
% For dealing with references we use the comment environment
\usepackage{verbatim}
\newenvironment{reference}{\comment}{\endcomment}
%\newenvironment{reference}{}{}
\newenvironment{slogan}{\comment}{\endcomment}
\newenvironment{history}{\comment}{\endcomment}

% For commutative diagrams we use Xy-pic
\usepackage[all]{xy}

% We use 2cell for 2-commutative diagrams.
\xyoption{2cell}
\UseAllTwocells

% We use multicol for the list of chapters between chapters
\usepackage{multicol}

% This is generally recommended for better output
\usepackage{lmodern}
\usepackage[T1]{fontenc}

% For cross-file-references

% Package for hypertext links:
\usepackage{hyperref}

% For any local file, say "hello.tex" you want to link to please

% Theorem environments.
%
\theoremstyle{plain}
\newtheorem{theorem}[subsection]{Theorem}
\newtheorem{proposition}[subsection]{Proposition}
\newtheorem{lemma}[subsection]{Lemma}

\theoremstyle{definition}
\newtheorem{definition}[subsection]{Definition}
\newtheorem{example}[subsection]{Example}
\newtheorem{exercise}[subsection]{Exercise}
\newtheorem{situation}[subsection]{Situation}

\theoremstyle{remark}
\newtheorem{remark}[subsection]{Remark}
\newtheorem{remarks}[subsection]{Remarks}

\numberwithin{equation}{subsection}

% Macros
%
\def\lim{\mathop{\mathrm{lim}}\nolimits}
\def\colim{\mathop{\mathrm{colim}}\nolimits}
\def\Spec{\mathop{\mathrm{Spec}}}
\def\Hom{\mathop{\mathrm{Hom}}\nolimits}
\def\Ext{\mathop{\mathrm{Ext}}\nolimits}
\def\SheafHom{\mathop{\mathcal{H}\!\mathit{om}}\nolimits}
\def\SheafExt{\mathop{\mathcal{E}\!\mathit{xt}}\nolimits}
\def\Sch{\mathit{Sch}}
\def\Mor{\mathop{\mathrm{Mor}}\nolimits}
\def\Ob{\mathop{\mathrm{Ob}}\nolimits}
\def\Sh{\mathop{\mathit{Sh}}\nolimits}
\def\NL{\mathop{N\!L}\nolimits}
\def\CH{\mathop{\mathrm{CH}}\nolimits}
\def\proetale{{pro\text{-}\acute{e}tale}}
\def\etale{{\acute{e}tale}}
\def\QCoh{\mathit{QCoh}}
\def\Ker{\mathop{\mathrm{Ker}}}
\def\Im{\mathop{\mathrm{Im}}}
\def\Coker{\mathop{\mathrm{Coker}}}
\def\Coim{\mathop{\mathrm{Coim}}}

% Boxtimes
%
\DeclareMathSymbol{\boxtimes}{\mathbin}{AMSa}{"02}

%
% Macros for moduli stacks/spaces
%
\def\QCohstack{\mathcal{QC}\!\mathit{oh}}
\def\Cohstack{\mathcal{C}\!\mathit{oh}}
\def\Spacesstack{\mathcal{S}\!\mathit{paces}}
\def\Quotfunctor{\mathrm{Quot}}
\def\Hilbfunctor{\mathrm{Hilb}}
\def\Curvesstack{\mathcal{C}\!\mathit{urves}}
\def\Polarizedstack{\mathcal{P}\!\mathit{olarized}}
\def\Complexesstack{\mathcal{C}\!\mathit{omplexes}}
% \Pic is the operator that assigns to X its picard group, usage \Pic(X)
% \Picardstack_{X/B} denotes the Picard stack of X over B
% \Picardfunctor_{X/B} denotes the Picard functor of X over B
\def\Pic{\mathop{\mathrm{Pic}}\nolimits}
\def\Picardstack{\mathcal{P}\!\mathit{ic}}
\def\Picardfunctor{\mathrm{Pic}}
\def\Deformationcategory{\mathcal{D}\!\mathit{ef}}

\setlength{\emergencystretch}{3em}
\begin{document}
\title[Stacks Project, Tag 07K7]{361 --- Unbounded right derived functors with finite cohomological dimension}
\maketitle
\noindent Copyright (C) 2005--2025 Johan de Jong. Stacks Project authors. Source: \href{https://stacks.math.columbia.edu/tag/07K7}{Tag 07K7}.

\noindent Editorial difficulty note (not author text). Difficulty H2: The proof constructs acyclic-term replacements for arbitrary unbounded complexes and proves the resulting functor preserves quasi-isomorphisms. Truncation cones and a hypercohomology spectral sequence isolate a finite cohomological window, after which separate distinguished-triangle arguments establish the sharp truncation and amplitude bounds. The unbounded case cannot be settled by ordinary bounded injective resolutions..

\noindent Editorial provenance note (not author text). History: excerpt from revision \texttt{a0444\allowbreak 6e57e\allowbreak c1fbc\allowbreak 252a8\allowbreak 71afc\allowbreak ec775\allowbreak 2fb28\allowbreak 07b14}, derived.tex, lines 9976--10100. Reference presentation was adapted; original-author context and core support blocks were retained or added. The mathematical wording is unchanged. Licensed under GNU FDL 1.2 or later; no invariant sections or cover texts. See ../licenses/COPYING. Context blocks reproduce author definitions and supporting results; they are not additional counted problems.

\begin{definition}
\label{definition-derived-functor}
In
Situation \ref{situation-classical}.
\begin{enumerate}
\item The {\it right derived functors of $F$} are the partial functors
$RF$ associated to cases (1) and (2) of
Situation \ref{situation-classical}.
\item The {\it left derived functors of $F$} are the partial functors
$LF$ associated to cases (3) and (4) of
Situation \ref{situation-classical}.
\item An object $A$ of $\mathcal{A}$ is said to be
{\it right acyclic for $F$}, or {\it acyclic for $RF$}
if $A[0]$ computes $RF$.
\item An object $A$ of $\mathcal{A}$ is said to be
{\it left acyclic for $F$}, or {\it acyclic for $LF$}
if $A[0]$ computes $LF$.
\end{enumerate}
\end{definition}

\begin{situation}
\label{situation-classical}
Here $F : \mathcal{A} \to \mathcal{B}$ is an additive functor between
abelian categories. This induces exact functors
$$
F : K(\mathcal{A}) \to K(\mathcal{B}), \quad
K^{+}(\mathcal{A}) \to K^{+}(\mathcal{B}), \quad
K^{-}(\mathcal{A}) \to K^{-}(\mathcal{B}).
$$
See Lemma \href{https://stacks.math.columbia.edu/tag/014X}{lemma additive exact homotopy category (Tag 014X)}.
We also denote $F$ the composition $K(\mathcal{A}) \to D(\mathcal{B})$,
$K^{+}(\mathcal{A}) \to D^{+}(\mathcal{B})$, and
$K^{-}(\mathcal{A}) \to D^-(\mathcal{B})$ of $F$ with the localization
functor $K(\mathcal{B}) \to D(\mathcal{B})$, etc. This situation leads
to four derived functors we will consider in the following.
\begin{enumerate}
\item The right derived functor of
$F : K(\mathcal{A}) \to D(\mathcal{B})$
relative to the multiplicative system $\text{Qis}(\mathcal{A})$.
\item The right derived functor of
$F : K^{+}(\mathcal{A}) \to D^{+}(\mathcal{B})$
relative to the multiplicative system $\text{Qis}^{+}(\mathcal{A})$.
\item The left derived functor of
$F : K(\mathcal{A}) \to D(\mathcal{B})$
relative to the multiplicative system $\text{Qis}(\mathcal{A})$.
\item The left derived functor of
$F : K^{-}(\mathcal{A}) \to D^{-}(\mathcal{B})$
relative to the multiplicative system $\text{Qis}^-(\mathcal{A})$.
\end{enumerate}
Each of these cases is an example of
Situation \ref{situation-derived-functor}.
\end{situation}

\begin{situation}
\label{situation-derived-functor}
Here $F : \mathcal{D} \to \mathcal{D}'$ is an exact functor of triangulated
categories and $S$ is a saturated multiplicative
system in $\mathcal{D}$ compatible with the structure
of triangulated category on $\mathcal{D}$.
\end{situation}

\begin{lemma}
\label{lemma-unbounded-right-derived}
Let $F : \mathcal{A} \to \mathcal{B}$ be a left exact functor of
abelian categories. Assume
\begin{enumerate}
\item every object of $\mathcal{A}$ is a subobject of an object
which is right acyclic for $F$,
\item there exists an integer $n \geq 0$ such that $R^nF = 0$,
\end{enumerate}
Then
\begin{enumerate}
\item $RF : D(\mathcal{A}) \to D(\mathcal{B})$ exists,
\item any complex consisting of right acyclic objects for $F$ computes $RF$,
\item any complex is the source of a quasi-isomorphism into a complex
consisting of right acyclic objects for $F$,
\item for $E \in D(\mathcal{A})$
\begin{enumerate}
\item $H^i(RF(\tau_{\leq a}E) \to H^i(RF(E))$ is an isomorphism
for $i \leq a$,
\item $H^i(RF(E)) \to H^i(RF(\tau_{\geq b - n + 1}E))$ is an isomorphism
for $i \geq b$,
\item if $H^i(E) = 0$ for $i \not \in [a, b]$ for some
$-\infty \leq a \leq b \leq \infty$, then $H^i(RF(E)) = 0$
for $i \not \in [a, b + n - 1]$.
\end{enumerate}
\end{enumerate}
\end{lemma}

\begin{proof}
Note that the first assumption implies that
$RF : D^+(\mathcal{A}) \to D^+(\mathcal{B})$ exists, see
Proposition \href{https://stacks.math.columbia.edu/tag/05TA}{proposition enough acyclics (Tag 05TA)}.
Let $A$ be an object of $\mathcal{A}$. Choose an injection $A \to A'$
with $A'$ acyclic. Then we see that $R^{n + 1}F(A) = R^nF(A'/A) = 0$ by
the long exact cohomology sequence. Hence we conclude that $R^{n + 1}F = 0$.
Continuing like this using induction we find that $R^mF = 0$ for all
$m \geq n$.

\medskip\noindent
We are going to use Lemma \href{https://stacks.math.columbia.edu/tag/07K6}{lemma replace resolution (Tag 07K6)} with the function
$d : \Ob(\mathcal{A}) \to \{0, 1, 2, \ldots \}$ given by
$d(A) = \max \{0\} \cup \{i \mid R^iF(A) \not = 0\}$.
The first assumption of Lemma \href{https://stacks.math.columbia.edu/tag/07K6}{lemma replace resolution (Tag 07K6)}
is our assumption (1). The second assumption of
Lemma \href{https://stacks.math.columbia.edu/tag/07K6}{lemma replace resolution (Tag 07K6)} follows from the fact
that $RF(A \oplus B) = RF(A) \oplus RF(B)$. The third assumption of
Lemma \href{https://stacks.math.columbia.edu/tag/07K6}{lemma replace resolution (Tag 07K6)} follows from the long exact
cohomology sequence. Hence for every complex $K^\bullet$ there exists a
quasi-isomorphism $K^\bullet \to L^\bullet$ into a complex of
objects right acyclic for $F$. This proves statement (3).

\medskip\noindent
We claim that if $L^\bullet \to M^\bullet$ is a quasi-isomorphism of
complexes of right acyclic objects for $F$, then
$F(L^\bullet) \to F(M^\bullet)$
is a quasi-isomorphism. If we prove this claim then we get statements
(1) and (2) of the lemma by
Lemma \href{https://stacks.math.columbia.edu/tag/06XN}{lemma find existence computes (Tag 06XN)}.
To prove the claim pick an integer $i \in \mathbf{Z}$.
Consider the distinguished triangle
$$
\sigma_{\geq i - n - 1}L^\bullet \to
\sigma_{\geq i - n - 1}M^\bullet \to Q^\bullet,
$$
i.e., let $Q^\bullet$ be the cone of the first map.
Note that $Q^\bullet$ is bounded below and that
$H^j(Q^\bullet)$ is zero except possibly for $j = i - n - 1$
or $j = i - n - 2$. We may apply $RF$ to $Q^\bullet$.
Using the second spectral sequence of
Lemma \href{https://stacks.math.columbia.edu/tag/015J}{lemma two ss complex functor (Tag 015J)}
and the assumed vanishing of cohomology (2) we conclude
that $H^j(RF(Q^\bullet))$ is zero except possibly for
$j \in \{i - n - 2, \ldots, i - 1\}$. Hence we see that
$RF(\sigma_{\geq i - n - 1}L^\bullet) \to RF(\sigma_{\geq i - n - 1}M^\bullet)$
induces an isomorphism of cohomology objects in degrees $\geq i$.
By Proposition \href{https://stacks.math.columbia.edu/tag/05TA}{proposition enough acyclics (Tag 05TA)} we know that
$RF(\sigma_{\geq i - n - 1}L^\bullet) = \sigma_{\geq i - n - 1}F(L^\bullet)$
and
$RF(\sigma_{\geq i - n - 1}M^\bullet) = \sigma_{\geq i - n - 1}F(M^\bullet)$.
We conclude that $F(L^\bullet) \to F(M^\bullet)$
is an isomorphism in degree $i$ as desired.

\medskip\noindent
Part (4)(a) follows from Lemma \href{https://stacks.math.columbia.edu/tag/05TC}{lemma negative vanishing (Tag 05TC)}.

\medskip\noindent
For part (4)(b) let $E$ be represented by the complex $L^\bullet$
of objects right acyclic for $F$. By part (2) $RF(E)$ is represented
by the complex $F(L^\bullet)$ and $RF(\sigma_{\geq c}L^\bullet)$
is represented by $\sigma_{\geq c}F(L^\bullet)$. Consider the
distinguished triangle
$$
H^{b - n}(L^\bullet)[n - b] \to
\tau_{\geq b - n}L^\bullet \to
\tau_{\geq b - n + 1}L^\bullet
$$
of Remark \href{https://stacks.math.columbia.edu/tag/08J5}{remark truncation distinguished triangle (Tag 08J5)}.
The vanishing established above gives that
$H^i(RF(\tau_{\geq b - n}L^\bullet))$ agrees with
$H^i(RF(\tau_{\geq b - n + 1}L^\bullet))$ for $i \geq b$.
Consider the short exact sequence of complexes
$$
0 \to
\Im(L^{b - n - 1} \to L^{b - n})[n - b] \to
\sigma_{\geq b - n}L^\bullet \to
\tau_{\geq b - n}L^\bullet \to 0
$$
Using the distinguished triangle associated to this
(see Section \href{https://stacks.math.columbia.edu/tag/014Z}{section canonical delta functor (Tag 014Z)})
and the vanishing as before we conclude that
$H^i(RF(\tau_{\geq b - n}L^\bullet))$ agrees with
$H^i(RF(\sigma_{\geq b - n}L^\bullet))$ for $i \geq b$.
Since the map $RF(\sigma_{\geq b - n}L^\bullet) \to RF(L^\bullet)$
is represented by $\sigma_{\geq b - n}F(L^\bullet) \to F(L^\bullet)$
we conclude that this in turn agrees with $H^i(RF(L^\bullet))$
for $i \geq b$ as desired.

\medskip\noindent
Proof of (4)(c). Under the assumption on $E$ we have
$\tau_{\leq a - 1}E = 0$ and we get the vanishing
of $H^i(RF(E))$ for $i \leq a - 1$ from part (4)(a).
Similarly, we have $\tau_{\geq b + 1}E = 0$ and hence
we get the vanishing of $H^i(RF(E))$ for $i \geq b + n$ from
part (4)(b).
\end{proof}
\end{document}
